\documentclass[11pt]{article}
\usepackage[margin=0.8in]{geometry}
\usepackage{times}
\usepackage[T1]{fontenc}
\usepackage[utf8]{inputenc}
\setlength{\parindent}{0pt}
\setlength{\parskip}{0.55em}
\pagestyle{empty}

\begin{document}

Department of Physics\\
Massachusetts Institute of Technology\\
Cambridge, MA 02139, USA

\today

Dear Editor,

We are pleased to submit a revised version of our Article, ``Interspecies Interactions Drive Community-Level Selection in Microbial Coalescence,'' for continued consideration in Nature Ecology \& Evolution. We thank the editor and reviewers for their careful evaluation of the manuscript. Their comments helped us clarify the conceptual framing, strengthen the analyses, and make the scope and limits of our interpretation more explicit.

We have revised the manuscript and Supplementary Information in response to each reviewer point, as detailed in the accompanying point-by-point response. The major revisions clarify the operational definition of community-level selection as origin-correlated persistence during coalescence, quantified by parental-community similarity and pairwise species-fate correlation, rather than as a claim about one exclusive biochemical mechanism. We also refined the interpretation of the nutrient perturbation experiments, now describing the nutrient gradient as perturbing net interaction intensity, invasion resistance, and environmental feedbacks, rather than treating it as a direct measurement of the generalized Lotka--Volterra interaction-strength parameter.

We also added or expanded analyses requested by the reviewers, including continuous similarity metrics and boundary-sensitivity analyses, event-matched additive null models, richness and norm-imbalance controls, pH-feedback modeling, dominant-species and pH controls, alternative interaction-matrix simulations allowing facilitative or reciprocal interaction structures, and taxonomic-distinctness analyses for natural sample-derived communities. We revised the figures and Supplementary Information to improve explanations of pairwise selection correlation, assembly filtering, invasion resistance, and the scope of the natural sample-derived community experiments.

These revisions preserve the central conclusion of the manuscript while making its assumptions and scope more precise: effective ecological coupling and environmental feedbacks can determine when microbial coalescence behaves as independent species sorting and when parental-community origin becomes predictive of persistence. The revised manuscript separates the phenomenological model parameter from experimentally measured invasion resistance, distinguishes operational community-level selection from specific mechanistic routes, and shows that the main transition is robust to several null models and model extensions.

The revised Article contains six main figures and no tables. The compiled revised manuscript is 29 pages, including Methods, references, and figure legends. The revised Supplementary Information is 54 pages and includes Extended Data Figs.~1--8 and Supplementary Figs.~1--44. A detailed point-by-point response to all reviewer comments is submitted as a separate response-to-reviewers file.

We confirm that this manuscript has not been published elsewhere and is not under consideration by another journal. No related manuscripts are under consideration or in press elsewhere. All authors have approved the revised manuscript and agree with its resubmission to Nature Ecology \& Evolution. We have no conflicts of interest to disclose. Jeff Gore is the corresponding author.

Thank you for considering our revised manuscript.

Sincerely,

Jeff Gore (Corresponding Author), Jinyeop Song, and Jiliang Hu\\
jeffgore@gmail.com

\end{document}
